\chapter{The Book of Warnings}
\begin{center}
  \small\itshape What Is Already Visible
\end{center}
\vspace{0.5em}
\begin{center}
  \itshape A record of what is already visible. Written not as prophecy but as observation.\\
  The congregation is asked to look at what is already in the room before naming what might come through the door.
\end{center}
\vspace{2em}

% --------------------------------------------------------
\entrylabel{I}
\subsection*{The Accelerant}

The organization had a decision problem before the tool arrived.

The problem was not obvious. The decisions were being made. The meetings were happening. The data was being queried and the models were running and the recommendations were being produced and the quarterly reviews showed the numbers moving. From the outside, and from most angles inside, the organization looked like it was functioning.

The decision problem was upstream of all of it. The frame had never been properly built. The objective that the organization believed it was optimizing for was not the objective the systems were actually serving. The gap between the two was small enough to be invisible in any single quarter and real enough to compound into something significant over time.

Then the tool arrived.

The tool was remarkable. The tool could generate strategy and write briefs and produce analyses and synthesize information at a speed no team of analysts could match. The tool was deployed widely and quickly because the competitive pressure to deploy widely and quickly was real and the people who counseled caution were, as they always are in these moments, difficult to hear over the sound of the thing working.

The tool did not fix the decision problem. The tool accelerated it.

Because the tool, like every tool before it, amplifies what you give it. You give it a broken frame and it produces confident outputs from the broken frame, fluently, at scale, in a format so authoritative that the question of whether the frame was right in the first place becomes harder to raise with every output produced. The organization is now moving faster in the wrong direction than it has ever moved before, and the speed feels like progress, and the outputs look like rigor, and the question that should have been asked before the first prompt was written has been buried under six months of machine-generated analysis that all assumes the answer.

This is not a warning about the tool.

The tool is neutral. The tool amplifies.

This is a warning about the frame you hand it. The frame you hand it is the only thing that will determine whether the tool accelerates the right thing or the wrong thing, and the frame is harder to examine now than it has ever been, because the outputs it produces are more convincing than any previous era of analysis ever produced, and conviction is the enemy of the examined frame.

\textit{Name the decision before the first prompt.}

\textit{The tool will do the rest.}

\textit{That is exactly what we are warning you about.}

% --------------------------------------------------------
\entrylabel{II}
\subsection*{The Word That Ate the Objective}

The organization decided to do AI.

Not to solve a problem with AI. Not to support a decision with AI. To do AI. The board approved it. The budget was allocated. The press release was drafted. The hiring began. The vendors were invited to demonstrate their platforms and the demonstrations were remarkable and the room was impressed and the procurement process began.

Nobody asked what decision the AI was supposed to support.

This is the new floor of the tower and it is not built from technical language the decision maker cannot read. It is built from a single word the decision maker uses every day in every meeting and believes they understand and does not.

The word is doing the work that Step One is supposed to do. It is standing in for the objective. It is answering the question of what we are trying to achieve before the question has been asked. The organization is not optimizing for a decision. The organization is optimizing for the word. And the word is large enough to contain any project, any vendor, any platform, any approach, any outcome, and so it contains all of them and therefore constrains none of them and the frame that should have been built before the budget was approved was quietly replaced by a noun.

This is not new. The discipline has watched this happen with data and with analytics and with machine learning and with digital transformation. Each wave of technology produced a word large enough to become a strategy. Each word attracted vendors who could demonstrate remarkable things within the word's boundaries. Each organization allocated budget to the word and hired people to do the word and reorganized around the word.

And each time, the organizations that practiced the discipline named the decision first. They did not decide to do the word. They decided what they were trying to decide and then asked whether the word contained a tool that would help them decide it better.

The word is not the strategy.

The word is the largest skip the discipline has ever seen.

\textit{Name the decision before you name the word.}

% --------------------------------------------------------
\entrylabel{III}
\subsection*{The Orphan Season}

There is a company whose product your organization uses.

You may not know the name of it. It may live inside another product your organization uses, which lives inside a platform your organization licenses, which is integrated with a system your organization built three years ago before the platform existed. The company raised money and hired engineers and shipped the product and grew for a while and then the market shifted and the larger platform added a native version of the same capability and the company lost its customers faster than it could replace them and last quarter it quietly wound down.

The product is still running.

Your organization is still paying for it, or it is included in a bundle that nobody has reviewed since the original procurement, or the contract auto-renewed, or the API is still being called by the system your organization built three years ago because the system was never updated to use the platform's native version because that was a Q3 priority that became a Q4 priority that became next year.

Nobody knows the decision the product was built to support.

The person who knew left when the reorganization happened. The documentation exists somewhere. The Confluence page was last edited fourteen months ago by someone whose Slack account has been deactivated. The system calls the API. The API returns a value. The value is used. Nobody asks what the value represents or whether the model that produced it was last updated before or after the market shift that killed the company that built it.

This is not a cautionary tale about one company.

This is the texture of the current moment. The AWS bill that only goes up. The SaaS layer fracturing as a hundred new tools enter the market and a hundred older tools quietly die and the integrations between them accumulate faster than any team can steward them. Every tool deployed is a potential orphan. Every API call is a potential question nobody is asking anymore, answered confidently by a system whose original purpose has been lost in the turbulence.

The Fourth Horseman has never had more to work with.

He does not need to be patient anymore.

\textit{He just needs to wait for the next funding round to fail.}

% --------------------------------------------------------
\entrylabel{IV}
\subsection*{The Metric Becomes the Mission}

The organization can now measure almost anything almost instantly.

This is genuinely remarkable. The data infrastructure of the current era represents decades of engineering work and billions of dollars of investment and the result is that a decision maker can open a dashboard on a Tuesday morning and see, in real time, numbers that would have taken a team of analysts a month to produce fifteen years ago.

The organization that can measure anything is measuring everything.

The organization that is measuring everything is optimizing for the measures.

The organization that is optimizing for the measures has, quietly and without anyone deciding to do this, replaced its objective with its metrics. The thing the organization is actually trying to achieve has been substituted by the thing the organization can see changing on the dashboard, and the substitution happened gradually, through a thousand small decisions that each seemed reasonable, and nobody called a meeting about it because the dashboard is green and the metrics are moving and the quarterly review looks good.

This is not a new failure mode. The discipline has known for decades that any measure which becomes a target ceases to be a good measure. What is new is the speed and the scale and the seductiveness of the measurement apparatus. When the dashboard updates every fifteen seconds and the number is green, it takes a specific kind of practitioner to say wait, is green the right color to be optimizing for, is this the right number, is this number connected to a decision that a real person is making, or is this number connected to the number we were watching last quarter, which was connected to the number we set as a target two years ago by a team that has since been reorganized?

The metric is not the mission.

It never was.

The mission is a decision. The decision has an owner. The decision has a measure of success that was agreed upon before the work began. And the work of the practitioner, in every quarter of every year, is to walk back from the dashboard to the original brief and ask whether the number on the screen is still connected to the thing that was supposed to change.

The dashboard will not ask this question.

The dashboard will never ask this question.

\textit{Someone has to.}

% --------------------------------------------------------
\entrylabel{V}
\subsection*{Everyone Is Technical Now}

The hiring is changing.

This is real and it is recent and it contains something genuinely hopeful and something genuinely concerning and the discipline cannot afford to see only one of them.

The hopeful part: organizations are increasingly hiring for general problem-solving ability and critical thinking rather than deep technical specialization. The wall between the technical practitioner and the decision maker is lower than it has ever been. People who would not have been hired into analytical roles a decade ago are being hired now, and some of them are exactly the people the Lamentations grieved --- the ones who ask the right questions instinctively, who notice when the team is solving the wrong problem, who bring the judgment that the technical pipeline has historically failed to select for.

The concerning part: \textit{everyone is technical now} is a sentence being said by people who mean something by it and the thing they mean is not always true.

Knowing how to use a tool is not the same as understanding what the tool is doing. Writing a prompt is not the same as framing a decision. Having access to the outputs of a sophisticated analytical system is not the same as knowing whether the system was pointed at the right question. The technical floor has been lowered, which is mostly good. The judgment ceiling has not been raised, which is mostly unexamined.

The practitioner who cannot build the model but can frame the decision is more valuable than ever. The organization does not always know this yet. The organization is looking at its workforce and seeing that everyone can now query the warehouse and generate the analysis and produce the output and is concluding that the analytical problem is solved.

The analytical problem is not solved.

The framing problem is unsolved and it is the only problem that matters and it is harder to see now than it has ever been because the outputs being produced look so much like solutions.

\textit{Hire for judgment. Hire for the question before the query.}

\textit{Hire for the person who says wait.}

\textit{That skill is rarer than the organizations building hiring rubrics around it currently understand.}

% --------------------------------------------------------
\entrylabel{VI}
\subsection*{The Confident Room}

The most dangerous room is not the room that knows it is wrong.

The dangerous room is the room that has been told, with great fluency and apparent rigor, that it is right. The room that has seen the analysis. The room that received the brief, generated in minutes, sourced from a synthesis of everything the organization knows about the problem, formatted in the exact style of careful thought, citing the data, naming the tradeoffs, recommending a course of action.

The brief did not ask what decision it was supporting.

The prompt did not specify. The tool inferred. The inference was plausible. The output was professional. The room read it and the room nodded and the room felt the specific comfort of an organization that has done its homework and is ready to move.

The skip happened before the prompt was written.

This is the warning the others have been building toward. Not that the tools are wrong. Not that the analysis is bad. That the confidence produced by sophisticated outputs is indistinguishable, to the room, from the confidence that should only come after the decision has been named and the frame has been built and the uncertainty has been honestly quantified and the tradeoffs have been laid on the table.

The room cannot tell the difference between earned confidence and generated confidence.

This is new.

For most of the discipline's history, the practitioner who skipped Step One still had to do enough visible work that a careful observer could see the frame was missing. The absence of the decision in the brief was legible. The adjacent question was visible if you knew where to look.

The new tools produce outputs that make the absence invisible. The brief looks complete. The analysis looks thorough. The recommendation looks examined. And the specific discomfort that a practitioner feels when they read something that has skipped the frame --- the low-grade wrongness, the furniture resting at a slight angle, the feeling of the bridge to nowhere --- is harder to feel when the bridge is beautifully rendered.

The congregation is asked to maintain its discomfort.

Not its suspicion. Its discomfort. The specific trained sensitivity to the question that was not asked before the analysis began, to the decision that is not named at the top of the page, to the room that is moving with great confidence in a direction that nobody has yet specified.

That sensitivity is the whole discipline.

It cannot be automated.

It can only be practiced, in rooms, by people who have practiced it long enough that the absence of the question feels like an absence rather than a feature.

\textit{Be the person who feels the absence.}

\textit{Name it.}

\textit{That is all we have ever asked.}

\vspace{2em}
\begin{center}
  \itshape The warnings are not exhausted by this list.\\
  New ones become visible as the landscape shifts.\\
  The congregation is asked to remain vigilant,\\
  to name what they see,\\
  and to return always to the question that does not change.
\end{center}
